\section{Introduction}
\label{sec:intro}

Machine learning has transitioned from an exploratory paradigm into the fundamental operational foundation of modern Network Intrusion Detection Systems (NIDS). Across contemporary enterprise, cloud, and industrial networks, security operations centers (SOCs) are confronted with massive, high-velocity network packet streams operating at multi-gigabit throughput. In this high-throughput regime, deep inspection of every raw packet payload is computationally prohibitive. Consequently, the industry standard relies on flow-based telemetry: statistical aggregation engines (such as Zeek, YAF, and CICFlowMeter \cite{sharafaldin2018toward}) summarize temporal sequences of bidirectional IP packets into compact, tabular flow records comprising features such as inter-arrival timings, packet length distributions, header byte totals, and protocol flags.

To classify these high-dimensional tabular records in real time, machine learning ensembles based on Gradient Boosted Decision Trees (GBDTs)---most prominently XGBoost \cite{chen2016xgboost}, LightGBM \cite{ke2017lightgbm}, and CatBoost \cite{prokhorenkova2018catboost}---have emerged as the undisputed state of the art. Across extensive benchmark evaluations on standardized intrusion datasets, including CSE-CIC-IDS2018 \cite{sharafaldin2018toward}, UNSW-NB15 \cite{moustafa2015unsw}, and CIDDS-001 \cite{ring2017flow, dang2021evaluating, dang2022studying}, tree-based ensembles consistently outperform deep neural network architectures (such as Multi-Layer Perceptrons, Convolutional Neural Networks, and TabNet), routinely establishing clean classification accuracy between $99.67\%$ and $99.87\%$ for GBDTs (and $99.53\%$ for Random Forest), with Area Under the ROC Curve (ROC-AUC) scores reaching up to $99.99\%$ (and $\ge 99.95\%$) on enterprise datasets like CSE-CIC-IDS2018 \cite{dang2022using, sharafaldin2018toward}, while achieving $98.42\%$--$99.25\%$ on hybrid and cloud environments (UNSW-NB15 \cite{moustafa2015unsw} and CIDDS-001 \cite{ring2017flow, dang2021evaluating, dang2022studying}).

\subsection{The Paradox of Near-Perfect Detection Accuracy}
Despite these near-perfect empirical evaluations, learning-based NIDS face an acute crisis of practical robustness when deployed against real-world adversaries. Foundational security studies and machine learning audits have long cautioned that near-perfect empirical detection scores on isolated benchmark datasets often fail to generalize to operational environments due to closed-world assumptions \cite{sommer2010outside} and shortcut learning on spurious features \cite{geirhos2020shortcut}. 

This discrepancy stems from three fundamental architectural and methodological phenomena:
\begin{enumerate}[label=(\roman*)]
    \item \textbf{Shortcut Learning on Testbed Artifacts:} Public benchmark datasets are overwhelmingly generated within synthetic testbed environments. In these environments, attack traces are generated using automated, scripted attack tools (such as LOIC, HOIC, Metasploit, and Hydra) that execute with rigid, deterministic transmission schedules. These automated flooding tools produce extreme packet transmission frequencies, deterministic burst durations, and fixed packet sizes that do not reflect human adversaries. As formalized by Geirhos et al. \cite{geirhos2020shortcut}, decision trees greedily partition feature space along these superficial volumetric shortcuts rather than learning invariant, semantic characteristics of the malicious exploits.
    \item \textbf{Data Leakage and Temporal Session Correlation:} Standard evaluation pipelines partition datasets using random train/validation/test sampling. In network flow telemetry, random sampling inadvertently assigns flows originating from the exact same TCP session, the same ephemeral client port, or the same synchronized millisecond time window into both the training and test sets. Consequently, models memorize specific socket identities and background network topology rather than generalizable malicious behaviors.
    \item \textbf{The Stationary Distribution Fallacy \& Cross-Environment Shift:} Conventional statistical learning assumes an independent and identically distributed (i.i.d.) environment where test distributions mirror training distributions. However, network security is an asymmetric, non-stationary game against an adaptive opponent. Furthermore, as demonstrated by Apruzzese et al.~\cite{apruzzese2022modeling, apruzzese2024evaluating} and Dang~\cite{dang2022detecting, dang2023learning, dang2025towards}, evaluations confined to single-testbed datasets fail when subjected to cross-environment shifts, revealing that models memorize testbed-specific artifacts rather than invariant intrusion semantics. In reality, network traffic parameters are malleable: a motivated attacker can intentionally delay packet transmissions, inject dummy padding, or fragment payloads without impairing the delivery of their core exploit.
\end{enumerate}

\subsection{The Adversary's Evasion Dilemma}
From an adversary's perspective, evading a deployed tree-based NIDS represents a unique, formidable challenge that cannot be resolved using standard adversarial machine learning techniques developed for computer vision or natural language processing. The attacker is constrained by two fundamental barriers:

\subsubsection{Non-Differentiability of Decision Tree Ensembles}
In continuous domains such as image classification, adversarial attacks rely on computing the exact gradient of the loss function with respect to the input features:
\begin{equation}
    \nabla_x \mathcal{L}(f(x), y)
\end{equation}
Algorithms such as the Fast Gradient Sign Method (FGSM) \cite{goodfellow2015explaining}, Projected Gradient Descent (PGD) \cite{madry2018towards}, and Carlini-Wagner (C\&W) \cite{carlini2017towards} utilize this gradient to perform gradient ascent toward the nearest classification boundary. 

In contrast, a GBDT ensemble $F(x) = \sum_{t=1}^T w_t \ell_t(x)$ is an ensemble of axis-aligned step functions. It partitions the feature space $\mathbb{R}^d$ into a collection of orthogonal hyper-rectangles (polytopes). Within the interior of each polytope, the ensemble prediction is constant, meaning the mathematical gradient is zero almost everywhere:
\begin{equation}
    \nabla_x F(x) = 0 \quad \text{for almost all } x \in \mathbb{R}^d
\end{equation}
At the polyhedral boundaries, the function exhibits discontinuous jump transitions where the classical gradient is undefined. Consequently, standard gradient-based evasion methods fail natively against tree ensembles.

\subsubsection{The Problem-Space Constraint Gap}
In network telemetry, tabular feature vectors $x \in \mathbb{R}^d$ do not exist in isolation; they are mathematical summaries computed by an aggregation function $\Phi$ applied to an underlying, physical sequence of network packets $P \in \mathcal{P}_{\text{valid}}$. 

As formalized in the seminal work by Pierazzi et al. \cite{pierazzi2020intrusion} and extended by Apruzzese et al. \cite{apruzzese2022modeling, apruzzese2023real}, an attacker operating in the problem space can only manipulate raw packets. Any valid adversarial feature perturbation $\Delta x$ must correspond to a physical packet transformation $P \to P^*$ that satisfies strict physical conservation laws:
\begin{itemize}
    \item \textbf{Protocol Immutability:} Attributes such as the Destination Port and Transport Protocol cannot be modified by the attacker if they seek to exploit a specific service (e.g., targeting an Apache web server on TCP port 80).
    \item \textbf{Directional Monotonicity:} In physical networks, an attacker can inject additional packets (dummy keep-alives) or insert padding bytes into application payloads. These actions can only \emph{increase} forward packet counts, forward byte totals, and total flow duration. An attacker cannot subtract bytes below the size of their exploit shellcode, nor can they transmit negative packets.
    \item \textbf{Non-Linear Conservation Laws:} Network flow features are tightly coupled through mathematical relationships. For instance:
    \begin{align}
        \text{Flow Bytes/s} &= \frac{\text{Total Fwd Bytes} + \text{Total Bwd Bytes}}{\text{Flow Duration} \times 10^{-6}} \label{eq:intro_byte_rate} \\
        \text{Flow Packets/s} &= \frac{\text{Total Fwd Packets} + \text{Total Bwd Packets}}{\text{Flow Duration} \times 10^{-6}} \label{eq:intro_pkt_rate} \\
        \text{Mean Fwd Packet Length} &= \frac{\text{Total Fwd Bytes}}{\text{Total Fwd Packets}} \label{eq:intro_mean_len}
    \end{align}
    If an adversarial algorithm perturbs Total Forward Bytes without simultaneously recalculating Flow Bytes/s and Mean Forward Packet Length, it generates a mathematically inconsistent ``phantom flow''. When evaluated by a real-world middlebox or stateful intrusion detector, these physically impossible flows are immediately discarded or flagged by rule-based protocol sanity checks.
\end{itemize}

\subsection{The Failure of Existing Paradigms}
Faced with these challenges, existing research in adversarial NIDS has remained trapped between two deeply flawed paradigms:
\begin{enumerate}
    \item \textbf{Unconstrained Feature-Space Optimization:} A large body of literature applies optimization algorithms directly in $\mathbb{R}^d$, including exact Mixed-Integer Linear Programming (MILP) solvers \cite{kantchelian2016evasion}, Zeroth-Order Optimization (ZOO) \cite{chen2017zoo}, and boundary-following search methods \cite{chen2020hopskipjump}. While exact MILP solvers operate in unconstrained feature space (yielding 0\% problem-space feasibility by design) and heuristic search algorithms achieve mathematical evasion on paper, they completely ignore protocol coupling. Their generated perturbations routinely violate basic physical conservation laws and monotonicity bounds, resulting in severe Physical Feasibility Rate (PFR) drops (measured at $17.8\%$--$25.0\%$ for benchmarked black-box attacks like ZOO and HopSkipJump, with unconstrained feature optimization yielding $0.0\%$ Real ASR) when tested against flow conservation checkers \cite{apruzzese2022modeling, pierazzi2020intrusion}.
    \item \textbf{Opaque Packet-Space Heuristics:} At the other extreme, packet-level mutation frameworks such as Traffic Manipulator \cite{han2021traffic} operate directly on packet captures (PCAPs) using Particle Swarm Optimization (PSO), while constraint-aware tabular search methods (such as CAA and MOEVA \cite{simonetto2024constrained, ghamizi2020moeva}) incorporate heuristic box constraints. While raw packet mutators guarantee routability, they treat the NIDS classifier as an opaque black box, requiring hundreds to thousands of queries ($1,240$--$1,850$ queries for evolutionary and swarm optimizers, and $680$ queries for HopSkipJump in our benchmark, taking up to several seconds per flow and minutes per batch) and failing to exploit the internal polyhedral partition geometry of tree ensembles.
\end{enumerate}

\subsection{Our Proposed Paradigm: Flow-Field Evasion (FFE)}
To bridge this fundamental divide between mathematical optimization and physical network reality, this article introduces \textbf{Flow-Field Evasion (FFE)}. 

Rather than perturbing tabular features arbitrarily or mutating packets blindly, FFE establishes a rigorous differential-geometric foundation for network flow evasion:
\begin{enumerate}
    \item We model the space of routable network flows through an ambient continuous relaxation $\widetilde{\mathcal{M}}_{\text{flow}} \subset \mathbb{R}^d$ constrained by network conservation laws, coupled with a discrete lattice projection $\mathcal{P}_{\mathcal{M}}$ onto valid packet-derived flow records $\mathcal{M}_{\text{flow}}$.
    \item We derive the closed-form \textbf{Infinitesimal Generators} (tangent vector fields) for four elementary physical network actions: Temporal Dilation ($\vec{v}_{\text{time}}$), Payload Padding ($\vec{v}_{\text{pad}}$), Chaff Packet Injection ($\vec{v}_{\text{chaff}}$), and TCP Window Morphing ($\vec{v}_{\text{win}}$).
    \item These generators define the \textbf{Physical Lie Cone} $\mathcal{K}_{\text{net}}(x) \subset T_x \widetilde{\mathcal{M}}_{\text{flow}}$. We prove that integrating non-negative combinations of these generators along $\mathcal{K}_{\text{net}}(x)$ preserves 100\% of physical conservation laws and protocol invariants by construction.
    \item By evaluating the directional Lie derivative $\mathcal{L}_{\vec{v}_k} F(x)$ of the tree ensemble's decision score along these physical generators, FFE steers attack flows along the optimal physical Lie-flow trajectory into the benign classification region with minimal query complexity.
\end{enumerate}

\subsection{Summary of Key Contributions}
The primary contributions of this work are summarized as follows:
\begin{itemize}
    \item \textbf{Physical Lie Flow Algebra:} We formalize the non-linear manifold of valid network flows $\mathcal{M}_{\text{flow}}$ and derive the exact closed-form infinitesimal generators for four elementary packet operations: Temporal Dilation, Payload Padding, Chaff Packet Injection, and TCP Window Morphing.
    \item \textbf{Universal Invariant Preservation Proof:} We prove mathematically that any trajectory generated by integrating non-negative linear combinations of these generators along the positive Lie Cone $\mathcal{K}_{\text{net}}(x)$ preserves 100\% of physical flow invariants, packet size bounds, and causal conservation laws by construction.
    \item \textbf{The Flow-Field Evasion (FFE) Algorithm:} We develop a directional Lie derivative steering algorithm that computes $\mathcal{L}_{\vec{v}_k} F(x)$ across the non-differentiable step boundaries of tree ensembles, guiding flows across the nearest polyhedral facets into the benign decision space with minimal operational transmission cost.
    \item \textbf{Multi-Dataset GPU Benchmark:} We implement and benchmark FFE on an NVIDIA RTX 5000 Ada GPU across three diverse benchmark corpora calibrated to empirical marginal distributions of enterprise (CSE-CIC-IDS2018), contemporary hybrid (UNSW-NB15), and cloud OpenStack (CIDDS-001/CIDDS-002) environments---against CUDA-accelerated XGBoost and CPU-multithreaded LightGBM models. FFE consistently achieves $100.0\%$ Attack Success Rate (ASR) and strictly $100.0\%$ Physical Feasibility Rate (PFR) in an average of only $6.8 - 17.8$ queries (including tree-split snapping evaluations) and under $70$ ms latency per flow.
    \item \textbf{Defensive Countermeasures:} We formulate two hardening countermeasures: Malleability-Penalized Tree Splitting (MPS) and Lie-Guided Adversarial Training (LGAT), demonstrating how defenders can mitigate flow-field steering down to an 11.0\% evasion rate (from 100\% under undefended models to 32.0\% under MPS, 18.0\% under LGAT, and 11.0\% under joint defense).
\end{itemize}

\subsection{Organization of the Paper}
The remainder of this paper is structured into five cohesive sections as follows:
\begin{itemize}
    \item \textbf{Section~\ref{sec:sota}: Related Work} surveys the landscape of learning-based NIDS, decision tree verification, problem-space constraints, and constraint-aware evasion benchmarks.
    \item \textbf{Section~\ref{sec:formulation}: Problem Formulation \& Problem-Space Semantics} formalizes the network flow feature schema, the physical flow manifold $\mathcal{M}_{\text{flow}}$, physical invariance constraints, and the adversary's threat model.
    \item \textbf{Section~\ref{sec:framework}: The Flow-Field Evasion (FFE) Framework} establishes the differential geometry of network flows, derives the infinitesimal Lie generators, establishes the Lie Cone of admissible actions, proves the Invariant Preservation Theorem, and details the algorithmic realization.
    \item \textbf{Section~\ref{sec:empirical_evaluation}: Empirical Benchmark, Ablation Studies \& Defenses} presents the GPU benchmark across CSE-CIC-IDS2018, UNSW-NB15, and CIDDS-001 against CUDA XGBoost and LightGBM models across five attack families (with CatBoost and Random Forest evaluated in transferability ablations), benchmarking FFE against five baselines, followed by ablations and defenses (MPS and LGAT).
    \item \textbf{Section~\ref{sec:conclusion}: Discussion, Real-World Deployment, and Conclusion} analyzes inline packet proxy translation, checksum synchronization, middlebox resilience, operational boundaries, ethical considerations, and concludes the paper.
\end{itemize}

